\scanpage{179}
\mspage{493}
% position: top continued paragraph
\mbox{s c e n c e}\quad (*)\quad \mbox{d e\quad l' a u t r e\quad t e r m e}, constitue un état de désé-
quilibre grave, et à la limite (quand ``l'absence'' en question est quasiment com-
plète, et prolongée) un état qui mène à la destruction de la chose (ou de \unsure{l'être|l'être)}
en laquelle ce déséquilibre a lieu, voire même de lui et de son entourage.

% position: upper middle paragraph
Ainsi, un état d' \mbox{a c t i o n} ininterrompue, qui n'alterne avec des
périodes suffisantes d' \mbox{i n a c t i o n}, de repos, mène à l'épuisement, la
maladie et (à la limite) à la mort - chose d'ailleurs qui a été des plus actuel-
les dernièrement, pour moi! (**) Mais inversement, un état d'inaction excessif
mène à un affaiblissement et à une sclérose des capacités et fonctions du corps
ou du psychisme (suivant les cas), et à la limite, à la destruction. Dans le
cas de mon ``incident-maladie'' d'ailleurs, j'ai un exemple simultané des \mbox{d e u \unsure{x}}
déséquilibres: action excessive de l'esprit, inaction du corps (et un repos suf-
fisant ni pour l'un ni pour l'autre...).

% position: middle lower paragraph
Cette ``explication'', dans ce cas d'espèce, de la ``philosophie'' équilibre-
déséquilibre du yin et du yang, reste superficielle, en ce sens qu'elle ne
touche pas à un parti-pris culturel invétéré, valorisant le terme\unsure{.} yang, l'action,
en l' \mbox{o p p o s a n t} au terme yin\inserted{,} l'inaction. Celle-ci est ressentie comme
une chose ``négative'', pas productive ni intéressante à aucun point de vue, ad-
mise à la rigueur comme un pis-aller, lequel s'impose hélas même à la meilleure
volonté du monde, puisqu'il faut quand-même se reposer de temps en temps pour
pouvoir continuer à s'investir dans l'action (sous peine justement, comme je
viens de l'expliquer, de surmenage et Dieu sait quoi encore...). En somme,
l'inaction est vue comme l'humble servante de l'action, indispensable hélas mais
à part çà indigne d'attention ni d'estime.
% position: footnote rule
\figure{\begin{tikzpicture}\draw(0,0)--(3,0);\end{tikzpicture}}
% position: first footnote
(*) (16 octobre) En fait, cette ``absence'' n'est jamais totale il me semble -
en aucune chose, le yin ni le yang n'est présent à l'état pur, sans la présence
simultanée de son complémentaire, si infime soit-elle. Le ``déséquilibre'' dont
je parle est donc caractérisé, non par l'absence totale d'un des deux termes
complémentaires (chose jamais réalisée), mais par un état de \mbox{f a i b l e s s e}
excessive de ce terme. Un autre type de déséquilibre, ou de morbidité, se présen-
te quand l'un \mbox{e t} l'autre terme sont ``absents'', ou plus précisément, sont
présents mais très faiblement. Ainsi, dans le cas du couple ``action-inaction'',
un état d' \mbox{a g i t a t i o n}, qui n'``agit'' pas à proprement parler (si ce \unsure{n'es|n'est}
pour se perpétuer lui-même, pour maintenir une confusion), tout en dispersant de
l'énergie, peut être considéré sans doute comme un tel déséquilibre ``par défaut''
(du yin \mbox{e t} du yang).

% position: second footnote
(**) Voir à ce sujet les deux premières notes (n$^\circ$ 98, 99) du Cortège XI, ``Le
défunt (toujours pas décédé...)''.
% position: lower left short mark
\margin{-}
% position: foot pencil
\margin{178}
